% Editable/compilable online at https://letx.app
\documentclass[10pt,a4paper]{article}

% --- Geometry & Margins ---
\usepackage[margin=0.75in, top=0.8in, bottom=0.8in]{geometry}

% --- Typography ---
\usepackage[T1]{fontenc}
\usepackage[utf8]{inputenc}
\usepackage{tgpagella} % Beautiful Palatino-based serif font
\usepackage{microtype} % Improved spacing and justification
\usepackage{setspace}
\setstretch{1.02} % Clean, compact line spacing for a 2-page document

% --- Colors ---
\usepackage{xcolor}
\definecolor{accent}{HTML}{1B365D} % Deep Navy Accent
\definecolor{highlight}{HTML}{2E4A62} % Muted Slate Blue
\definecolor{textgray}{HTML}{4A4A4A} % Muted text color for sub-info

% --- Section Styling ---
\usepackage{titlesec}
\titleformat{\section}
  {\color{accent}\normalfont\Large\bfseries}
  {\thesection}{1em}{}
  [\vspace{-0.4em}\color{accent}\titlerule]

\titleformat{\subsection}
  {\color{accent}\normalfont\large\bfseries}
  {\thesubsection}{0.8em}{}

\titlespacing*{\section}{0pt}{1.0em}{0.5em}
\titlespacing*{\subsection}{0pt}{0.8em}{0.4em}

% --- Callout Box for Executive Summary ---
\usepackage[most]{tcolorbox}
\newtcolorbox{summarybox}{
  colback=accent!5,
  colframe=accent,
  arc=0mm,
  outer arc=0mm,
  top=2.5mm,
  bottom=2.5mm,
  left=3.5mm,
  right=3.5mm,
  boxrule=0pt,
  leftrule=3.5pt,
  before=\par\noindent\vspace{0.4em},
  after=\vspace{0.4em}
}

% --- Lists and Spacing ---
\usepackage{enumitem}
\setlist[itemize]{leftmargin=1.5em, itemsep=0.2em, topsep=0.15em}
\setlist[enumerate]{leftmargin=1.5em, itemsep=0.2em, topsep=0.15em}

% --- Hyperlinks ---
\usepackage{hyperref}
\hypersetup{
  colorlinks=true,
  linkcolor=accent,
  citecolor=accent,
  urlcolor=accent
}

% --- Footnotes and Headers ---
\usepackage{fancyhdr}
\pagestyle{fancy}
\fancyhf{}
\renewcommand{\headrulewidth}{0pt}
\fancyfoot[C]{\color{textgray}\small Alex J. Mercer --- Research Statement --- Page \thepage}

\begin{document}

% --- Title Block ---
\begin{center}
  {\fontsize{24pt}{28pt}\selectfont\color{accent}\bfseries Dr. Alex J. Mercer} \\
  \vspace{0.2em}
  {\fontsize{13pt}{16pt}\selectfont\scshape Research Statement} \\
  \vspace{0.4em}
  {\small \color{textgray} 
    Department of Computer Science $\cdot$ Stanford University \\
    \href{mailto:amercer@stanford.edu}{amercer@stanford.edu} $\cdot$ \href{https://alexmercer.github.io}{alexmercer.github.io} $\cdot$ +1 (650) 555-0199
  } \\
  \vspace{0.6em}
  \color{accent}\hrule height 1.2pt
\end{center}

\vspace{0.2em}

% --- Executive Summary Callout ---
\begin{summarybox}
  \noindent\textbf{Research Vision:} My research lies at the intersection of deep learning theory and robust systems engineering. Specifically, I design algorithms that bridge the gap between theoretical guarantees in machine learning and the unpredictable reality of open-world deployments. Over the next five years, my goal is to build \textbf{provably safe, resource-efficient, and self-adaptive autonomous systems} that can safely interact with dynamic environments.
\end{summarybox}

% --- Section: Overview ---
\section{Overview}
Modern artificial intelligence has achieved remarkable success in controlled laboratory settings. However, when deployed in high-stakes, real-world environments---such as autonomous driving, healthcare, and robotics---these systems frequently fail due to distribution shifts, adversarial inputs, and hardware constraints. 

My research program addresses these vulnerabilities through a two-pronged approach:
\begin{enumerate}
  \item \textbf{Robustness and Safety Foundations:} Developing mathematical frameworks that provide formal out-of-distribution (OOD) generalization guarantees.
  \item \textbf{Systems-Aware Deep Learning:} Co-designing neural architectures and training algorithms that optimize latency and energy efficiency on edge hardware without sacrificing performance.
\end{enumerate}
To date, my research has led to 8 publications in top-tier venues (NeurIPS, ICML, CVPR, and ICLR), including one Outstanding Paper Award. My open-source libraries have been adopted by over 50 academic labs and industrial research teams.

% --- Section: Past & Current Research ---
\section{Past \& Current Research}
My work is characterized by bridging theoretical machine learning principles with practical systems engineering. 

\subsection{Robustness under Distribution Shift}
Standard machine learning assumes that training and testing data are drawn from the same distribution. In practice, this assumption is regularly violated. To tackle this, I developed a framework for \emph{Invariant Feature Attribution} [1]. By constraining neural networks to learn causal features rather than spurious correlations, we improved out-of-distribution accuracy by $14\%$ on safety-critical computer vision benchmarks. This work mathematically formalizes the connection between gradient-based attribution methods and causal invariance, providing the first generalization bounds under covariate shift.

\subsection{Computational Efficiency for Vision-Language Models}
Large-scale vision-language models (VLMs) offer rich representations but are computationally prohibitive for edge devices. In [2], I introduced \emph{SparseFlow}, a dynamic structural pruning algorithm that compresses VLMs during runtime. Unlike static pruning, which permanently removes parameters, SparseFlow selectively activates sub-networks depending on the input complexity. This achieves a $3.5\times$ reduction in computational latency and a $60\%$ reduction in memory footprint, enabling real-time deployment on standard mobile processors while retaining $98.5\%$ of the original model's accuracy.

% --- Section: Research Agenda ---
\section{Research Agenda \& Future Directions}
As a faculty member, I will build an interdisciplinary research group focusing on the lifecycle of trustworthy machine learning, spanning from mathematical proofs to deployment.

\subsection{Short-Term Focus (Years 1--3): Provably Safe Foundation Models}
As foundation models are increasingly integrated into interactive environments (e.g., LLM agents in robotics), ensuring safety is paramount. I plan to develop \emph{Neural Guardrails}---a framework that wraps around pre-trained foundation models and monitors their outputs in real-time. By utilizing conformal prediction and runtime verification, these guardrails will guarantee, with user-defined probability, that the system's actions remain within a safe state space. This project will be supported by collaborations with roboticists and control theorists to translate safety bounds to physical actuators.

\subsection{Long-Term Vision (Years 4+): Autonomic Self-Healing ML Systems}
The ultimate goal of my research is to create machine learning systems that can continuously adapt to new environments without human intervention. Standard fine-tuning often leads to catastrophic forgetting. I will investigate methods for \emph{Continual Modular Learning}, where the model dynamically allocates new sub-networks to represent new concepts while freezing existing modules. By formulating this as a Bayesian task-routing problem, we aim to build systems that automatically detect when their performance degrades, isolate the failure mode, and patch their own weights using streaming data.

% --- Section: Broader Impacts ---
\section{Broader Impacts}
\subsection{Education, Mentorship, and Diversity}
I am deeply committed to broadening participation in computer science. Throughout my PhD and postdoc, I have mentored 5 undergraduate students (including 3 from underrepresented backgrounds), resulting in two co-authored workshop papers. As a faculty member, I will design a new curriculum on \emph{Trustworthy Systems-Aware Machine Learning} to teach students how to balance algorithmic efficiency with societal safety. I will also partner with local high schools to run summer workshops introducing AI ethics and coding to underrepresented groups.

\subsection{Open-Source Software and Industrial Collaboration}
Democratizing research is key to accelerating scientific progress. I actively maintain the \texttt{RobustTorch} library, which simplifies the evaluation of model robustness under distribution shift. I will continue to foster collaborations with industry partners (having previously worked with researchers at Google DeepMind and NVIDIA) to ensure that the theoretical frameworks developed in my lab address real-world engineering bottlenecks and are quickly translated into production.

% --- Section: Selected Publications ---
\section{Selected Publications}
\begin{enumerate}[label={\textbf{[\arabic*]}}, leftmargin=*, itemsep=0.4em]
  \item \textbf{Alex J. Mercer}, Yu-An Chen, and Li-Min Tseng. ``Provably Robust Foundation Models via Neural Guardrails.'' \emph{International Conference on Machine Learning (ICML)}, 2025. \\
  \textit{Introduced a conformal-prediction framework to guarantee safety bounds in foundation models.} \hfill \color{accent}\textbf{[Outstanding Paper Award]}
  
  \item \textbf{Alex J. Mercer} and Sarah Jenkins. ``SparseFlow: Dynamic Pruning for Vision-Language Models.'' \emph{IEEE Conference on Computer Vision and Pattern Recognition (CVPR)}, 2024. \\
  \textit{Designed a dynamic sub-network activation scheme for edge computing, achieving $3.5\times$ latency reduction.}
  
  \item \textbf{Alex J. Mercer}, Marcus Vance, and David K. Harrison. ``Invariant Feature Attribution for Causal Out-of-Distribution Generalization.'' \emph{Advances in Neural Information Processing Systems (NeurIPS)}, 2023. \\
  \textit{Proposed causal invariance constraints to improve neural network generalization under covariate shifts.}
  
  \item \textbf{Alex J. Mercer} and David K. Harrison. ``Quantifying and Mitigating Spurious Correlations in Deep Classifiers.'' \emph{International Conference on Learning Representations (ICLR)}, 2022. \\
  \textit{Identified structural vulnerabilities in deep vision classifiers and proposed adversarial training mitigations.}
\end{enumerate}

\end{document}
